% Numerical continuation proof for the frozen IID verifier.
% Intended to follow the exact-arithmetic pruning and monotonicity lemmas.
\subsection{Certified evaluation of the continuation terms}
\label{sec:certificate-continuation-roundoff}

This subsection concerns the frozen instance with $N_t=640$, state knots
that are integer multiples of $2^{-30}$, and finite rate endpoints at most
$8192$. It justifies the continuation calculations in
\texttt{verify.cpp}, including the tests that return a zero continuation.
Recall the rounding allowances
\[
 \varepsilon_{\mathrm{in}}=2^{-48},\qquad
 \varepsilon_{\mathrm{ar}}=2^{-32}
       =65536\varepsilon_{\mathrm{in}},\qquad
 \mathfrak u=2^{-53}=\varepsilon_{\mathrm{in}}/32.
\]
The code calls the first two quantities \texttt{eta} and \texttt{eps}.
All floating-point operations below are binary64 operations in
round-to-nearest mode with gradual underflow (no flush-to-zero or
denormals-are-zero mode), without contraction or reassociation. We write
$\operatorname{fl}$ for the result of one such operation. The usual
bound, also valid when its result is subnormal, is
\begin{equation}
 |\operatorname{fl}(z)-z|
 \le \mathfrak u|z|+2^{-1075}.
 \label{eq:certificate-fp-model}
\end{equation}
None of the operations considered below overflows. The displayed bounds
retain ample strict slack for the additive term in
\eqref{eq:certificate-fp-model}; in particular, a positive state coordinate
satisfies $y\ge 2^{-30}$, and
$2^{-1075}<2^{-990}\varepsilon_{\mathrm{in}}y$.

\paragraph{Exact quantities and validated inputs.}
Put $\Delta=1/640$ and $t=k\Delta$, where $0\le k<640$.
At a nonterminal state, $0\le x<y\le1$.
The divisions that form $x$ and $y$ from the integer knots and $2^{30}$
are exact. Write $b$ for the exact drop rate on this time interval and
$D=1-\theta(t)$. The input validator establishes
\[
 \tfrac14<b<2,\qquad \tfrac{17}{20}<c<\tfrac9{10},\qquad
 0\le D\le1.
\]
For a finite exact endpoint $r\in[0,8192]$, define
\[
 E(r)=e^{-\Delta r},\quad
 I_r=\int_0^\Delta e^{-rs}\,ds.
\]
Thus $0<E(r)\le1$, $0<I_r\le\Delta$, and $0\le (r+e^{-r}-1)\le8192$.
The frozen values of $b,c,D,r,E(r),I_r,(r+e^{-r}-1)$ differ from these exact
quantities by less than $\varepsilon_{\mathrm{in}}$ in absolute value.
The source entry \texttt{vv[h]} is the frozen approximation of
$r+e^{-r}-1$. Both arrays \texttt{bs} and \texttt{bl} approximate the same exact $b$.
The frozen $\widehat\Delta$ satisfies the stronger relative bound
\[
 |\widehat\Delta-\Delta|<\Delta\varepsilon_{\mathrm{in}}.
\]
The computed $\widehat t=\operatorname{fl}(k\widehat\Delta)$ consequently
satisfies
\begin{equation}
 |\widehat t-t|
 \le k\Delta\varepsilon_{\mathrm{in}}
   +\mathfrak u k\Delta(1+\varepsilon_{\mathrm{in}})+2^{-1075}
 <2\varepsilon_{\mathrm{in}}.
 \label{eq:certificate-time-error}
\end{equation}
Input validation uses rational enclosures; these input-error assumptions
therefore do not assume accuracy of a platform exponential routine.

Recall the exact continuation quantities
\begin{align*}
 \varphi(r)&=1-y+y\bigl[t+I_r+E(r)(1-t-\Delta)\bigr],\\
 y(r)&=\frac{yE(r)}{\varphi(r)},\qquad
 \rho(r)=\frac{x}{y}+\frac{b\Delta (r+e^{-r}-1)}c,\\
 \widetilde x(r)&=y(r)\rho(r),\qquad
 x(r)=\min\{y(r),\widetilde x(r)\}.
\end{align*}
We will use
\begin{equation}
 E(r)\le\varphi(r)\le1,
 \qquad 0\le y(r)\le1.
 \label{eq:certificate-phi-range}
\end{equation}
For the lower inequality, subtract $E(r)$ from $\varphi(r)$ and use
$I_r\ge\Delta E(r)$; for the upper inequality use
$I_r\le\Delta$ and $t+\Delta\le1$.

\paragraph{Elementary operations used by the guards.}
Here and below an intermediate is evaluated in the order prescribed by
the source. Repeated application of
\eqref{eq:certificate-fp-model} gives the following bounds. The table
also records sufficient estimates for each operation, so that no
unquantified error-propagation assertion is needed. Error bounds in its
last column are multiples of $\varepsilon_{\mathrm{in}}$.
\begin{center}
\begin{tabular}{lll}
\hline
Computed intermediate & Exact target & Error bound\\
\hline
$\operatorname{fl}(\widehat b\widehat\Delta)$
 &$b\Delta$ &$4\Delta$\\
$\operatorname{fl}(\operatorname{fl}(\widehat b\widehat\Delta)\texttt{vv[h]})$
 &$b\Delta (r+e^{-r}-1)$ &$54$\\
$\operatorname{fl}(y\widehat b)$ &$yb$ &$9y/8$\\
$\operatorname{fl}(\operatorname{fl}(y\widehat b)\widehat\Delta)$
 &$yb\Delta$ &$4\Delta y$\\
$\operatorname{fl}(\operatorname{fl}(\operatorname{fl}(y\widehat b)
       \widehat\Delta)\texttt{vv[h]})$
 &$yb\Delta (r+e^{-r}-1)$ &$54y$\\
$\operatorname{fl}(\widehat c x)$ &$cx$ &$2y$\\
$\operatorname{fl}(\widehat D-\widehat b)$ &$D-b$ &$3$\\
$\operatorname{fl}(y\operatorname{fl}(\widehat D-\widehat b))$
 &$y(D-b)$ &$4y$\\
\hline
\end{tabular}
\end{center}
For example, the first line follows from
\[
 \Delta\varepsilon_{\mathrm{in}}
 +(2+\varepsilon_{\mathrm{in}})\Delta\varepsilon_{\mathrm{in}}
 +\mathfrak u(2+\varepsilon_{\mathrm{in}})
       \Delta(1+\varepsilon_{\mathrm{in}})+2^{-1075}
 <4\Delta\varepsilon_{\mathrm{in}}.
\]
For the second line the propagated input error and the new rounding
error are at most
\[
 4\Delta\,8192\varepsilon_{\mathrm{in}}
 +0.004\varepsilon_{\mathrm{in}}
 +25.61\mathfrak u+2^{-1075}
 <54\varepsilon_{\mathrm{in}}.
\]
The same calculation multiplied by $y$ proves the fifth line.
The other lines use $|D-b|\le2$, $c<0.9$, and
$\mathfrak u=\varepsilon_{\mathrm{in}}/32$.
The final additions in the code therefore give
\begin{align}
 |\widehat p-p|&<8\varepsilon_{\mathrm{in}}y,
 &p&=cx+y(D-b), \label{eq:certificate-p-error}\\
 |\widehat n(r)-n(r)|&<64\varepsilon_{\mathrm{in}}y,
 &n(r)&=cx+yb\Delta (r+e^{-r}-1), \label{eq:certificate-n-error}
\end{align}
because $|p|\le3y$ and $0\le n(r)\le26.5y$.
The right-hand side of the leaf clipping test is
$\widehat n_{\mathrm{clip}}=\operatorname{fl}(\operatorname{fl}(\widehat c y)
                    +\varepsilon_{\mathrm{ar}}y)$.
Multiplication by $\varepsilon_{\mathrm{ar}}$ is exact, and
\begin{equation}
 |\widehat n_{\mathrm{clip}}-(cy+\varepsilon_{\mathrm{ar}}y)|
 <2\varepsilon_{\mathrm{in}}y.
 \label{eq:certificate-clip-threshold-error}
\end{equation}
The factor $y$ in both bounds is essential for small state coordinates.

For the node guard the code computes
$\widehat\rho=\operatorname{fl}(\operatorname{fl}(x/y)
 +\operatorname{fl}(\widehat\omega/\widehat c))$,
where $\widehat\omega$ is the second intermediate in the table.
Since $b\Delta (r+e^{-r}-1)\le25.6$, $(b\Delta (r+e^{-r}-1))/c<31$, and
$\widehat c>0.84$, division has error at most
\[
 \frac{(54+31)\varepsilon_{\mathrm{in}}}{0.84}
 +32\mathfrak u+2^{-1075}.
\]
Including the division $x/y$ and the final addition yields
\begin{equation}
 |\widehat\rho-\rho(r)|<128\varepsilon_{\mathrm{in}}.
 \label{eq:certificate-ratio-error}
\end{equation}
Likewise, for the product in the upper-endpoint guard,
\begin{equation}
 |\operatorname{fl}(\widehat r\widehat\Delta)-r\Delta|
 <\bigl(0.002+12.8+0.401\bigr)\varepsilon_{\mathrm{in}}
 <16\varepsilon_{\mathrm{in}}.
 \label{eq:certificate-rate-product-error}
\end{equation}

\paragraph{Validity of the zero-continuation and range tests.}
If the leaf test skips interpolation, then $\widehat n(r)\ge\widehat n_{\mathrm{clip}}$.
Equations \eqref{eq:certificate-n-error} and
\eqref{eq:certificate-clip-threshold-error} imply
\[
 n(r)-cy\ge
  (\varepsilon_{\mathrm{ar}}-66\varepsilon_{\mathrm{in}})y>0.
\]
Thus $\rho(r)\ge1$, the exact successor is on the diagonal, and its
continuation is zero. If the leaf test instead evaluates the successor,
the reverse estimates give
\[
 \rho(r)<1+
 \frac{\varepsilon_{\mathrm{ar}}+66\varepsilon_{\mathrm{in}}}{c}
 <1+2\varepsilon_{\mathrm{ar}}.
\]
For the node test, $\widehat\rho\ge1+\varepsilon_{\mathrm{ar}}$ implies
$\rho(r)>1$ by \eqref{eq:certificate-ratio-error}, so all its descendant
continuations are zero by the monotonicity of $\rho$.
Otherwise $\rho(r)<1+2\varepsilon_{\mathrm{ar}}$ again.
In every case where a lower-endpoint successor is evaluated,
$(r+e^{-r}-1)\ge r-1$ now implies
\begin{equation}
 r\Delta\le\Delta+\frac cb\rho(r)
 <\frac1{640}+3.6(1+2\varepsilon_{\mathrm{ar}})<4.
 \label{eq:certificate-successor-rate-range}
\end{equation}
For an upper endpoint, the code evaluates its successor only if
$\operatorname{fl}(\widehat r\widehat\Delta)
 \le4-\varepsilon_{\mathrm{ar}}$.
Equation \eqref{eq:certificate-rate-product-error} proves $r\Delta<4$
in this case as well. Otherwise the node calculation uses the valid
lower bound zero for that successor's $y$-coordinate.

\paragraph{Errors in $\varphi$ and the successor coordinates.}
The following table follows the parenthesization of
\texttt{ph}. Its error bounds are again in units of
$\varepsilon_{\mathrm{in}}$.
\begin{center}
\begin{tabular}{lll}
\hline
Computed intermediate & Exact target & Error bound\\
\hline
$\operatorname{fl}(\operatorname{fl}(1-\widehat t)-\widehat\Delta)$
 &$1-t-\Delta$ &$3$\\
$\operatorname{fl}(\widehat E\widehat K)$
 &$E(1-t-\Delta)$ &$5$\\
$\operatorname{fl}(\widehat t+\widehat I)$ &$t+I_r$ &$4$\\
$\operatorname{fl}(\widehat T+\widehat E\widehat K)$
 &$t+I_r+E(1-t-\Delta)$ &$10$\\
$\operatorname{fl}(y\widehat Z)$ &$y[t+I_r+E(1-t-\Delta)]$ &$11y$\\
\hline
\end{tabular}
\end{center}
In the fourth line $\widehat E\widehat K$ denotes the already rounded
second-line product. Here $\widehat K,\widehat T,\widehat Z$ denote
the corresponding previously computed quantities.
For completeness, the first-line estimate is bounded by
$2+\Delta+(1+\Delta)/32$ in the stated units, with an arbitrarily
small allowance for $\varepsilon_{\mathrm{in}}$ and the subnormal term;
the next lines are bounded respectively by
$3+1+1/32$, $2+1+1/32$, $5+4+1/32$, and
$10y+y/32$, with the same strict slack.
The quantity $1-y$ is an exactly representable multiple of $2^{-30}$.
The last addition thus gives
\begin{equation}
 |\widehat\varphi(r)-\varphi(r)|
 <16\varepsilon_{\mathrm{in}}.
 \label{eq:certificate-phi-error}
\end{equation}

For every evaluated successor, \eqref{eq:certificate-successor-rate-range}
and \eqref{eq:certificate-phi-range} give
$\varphi(r)>e^{-4}>1/64$, hence
$\widehat\varphi(r)>1/65$.
The rounded numerator $\widehat a=\operatorname{fl}(y\widehat E)$
satisfies $|\widehat a-yE|<2\varepsilon_{\mathrm{in}}y$.
Using $y(r)\le1$, the error before rounding the division is at most
\[
 \left|\frac{\widehat a}{\widehat\varphi}-y(r)\right|
 \le\frac{|\widehat a-yE|+y(r)|\widehat\varphi-\varphi|}
             {\widehat\varphi}
 <65(2y+16)\varepsilon_{\mathrm{in}}.
\]
The division adds less than $2\mathfrak u+2^{-1075}$, proving
\begin{equation}
 |\widehat y(r)-y(r)|<2048\varepsilon_{\mathrm{in}}.
 \label{eq:certificate-successor-y-error}
\end{equation}

The leaf computes an \emph{unclipped} $x$-coordinate by dividing
$\operatorname{fl}(\widehat E\widehat n)$ by
$\operatorname{fl}(\widehat c\widehat\varphi)$.
Under its evaluation guard, $n=cy\rho<y$ and
$\widetilde x(r)\le1+2\varepsilon_{\mathrm{ar}}<1.001$.
Equation \eqref{eq:certificate-n-error} and one multiplication give
numerator error less than $66\varepsilon_{\mathrm{in}}y$.
Equation \eqref{eq:certificate-phi-error} gives denominator error less
than $16\varepsilon_{\mathrm{in}}$, and the computed denominator is
greater than $1/80$. The same quotient identity therefore bounds the
error before the final rounding by
\[
 80\bigl(66y+16\cdot1.001\bigr)\varepsilon_{\mathrm{in}}.
\]
Including the division rounding yields
\begin{equation}
 |\widehat x(r)-\widetilde x(r)|
 <8192\varepsilon_{\mathrm{in}}.
 \label{eq:certificate-successor-x-error}
\end{equation}

For a finite upper endpoint that passes its range guard, the node
instead computes
$\widehat x_{\mathrm{low}}
 =\operatorname{fl}(\widehat y(r_+)\min\{1,\widehat\rho(r_-)\})$.
The map $z\mapsto\min\{1,z\}$ is $1$-Lipschitz.
Combining \eqref{eq:certificate-ratio-error} and
\eqref{eq:certificate-successor-y-error}, then allowing one product
rounding, gives
\begin{equation}
 \bigl|\widehat x_{\mathrm{low}}
    -y(r_+)\min\{1,\rho(r_-)\}\bigr|
 <4096\varepsilon_{\mathrm{in}}.
 \label{eq:certificate-node-x-error}
\end{equation}
For example, the coefficient of $\varepsilon_{\mathrm{in}}$ is bounded
by $2048(1+128\varepsilon_{\mathrm{in}})+128+2/32$, which is below
$4096$. When the upper-endpoint guard is not used, the assigned value
is exactly zero and needs no floating-point estimate.

\paragraph{Shifted interpolation and multiplication.}
Let $\widetilde v_{k+1}$ be the ordinary triangular interpolant already
defined above; its vertex values are stored integers divided by $2^{40}$.
The monotonicity passes ensure $\widetilde v_{k+1}\le0$, that
$\widetilde v_{k+1}$ is nonincreasing in its first coordinate and
nondecreasing in its second coordinate, and that it is zero on the diagonal.
The interpolation routine rounds $\widehat y+\varepsilon_{\mathrm{ar}}$
up to a multiple of $2^{-30}$ and
$\widehat x-\varepsilon_{\mathrm{ar}}$ down to such a multiple,
then clips the former to $[0,1]$ and the latter below at zero.
The additions incur less than
$\varepsilon_{\mathrm{in}}/8$ error; multiplication by $2^{30}$ is
exact. Equations \eqref{eq:certificate-successor-y-error}--
\eqref{eq:certificate-node-x-error} and
$\varepsilon_{\mathrm{ar}}=65536\varepsilon_{\mathrm{in}}$
therefore put the integer query above the exact required
$y$-coordinate and below the exact required $x$-coordinate.
Clipping at $0$ or $1$ preserves these inequalities.

For a leaf with $\rho(r)<1$, the required $x$-coordinate is
$\widetilde x(r)=x(r)\le y(r)$, so the queried value is at most
$\widetilde v_{k+1}(y(r),x(r))$. If the routine returns zero because its integer
$x$-coordinate is no smaller than its integer $y$-coordinate, these
coordinate inequalities force the exact successor to be on the
diagonal; the zero return is then valid. For a leaf with $\rho(r)\ge1$,
the true continuation is already zero. In this case either a zero
return or any nonpositive interpolated value is conservative, even
though the code formed the raw coordinate $\widetilde x(r)$ rather
than its clipped version.
The identical argument applies to a node when $\rho(r_-)<1$ and the
required coordinate is $y(r_+)\rho(r_-)$; when $\rho(r_-)\ge1$, every
descendant continuation is zero. These are exactly the cases used by
the exact-arithmetic node bound.

Inside a containing triangle, the code forms nonnegative integer
barycentric numerators with their exact sum equal to the positive
denominator, and takes a downward integer quotient. Thus the returned
integer $v$ satisfies $v/2^{40}\le\widetilde v_{k+1}$ at the integer query. Integer
range bounds are established separately below.

Finally, the code uses
\[
 M=\left\lceil
     \operatorname{fl}\bigl((\widehat\varphi+
                       \varepsilon_{\mathrm{ar}})2^{40}\bigr)
    \right\rceil+1,
\]
where the parenthesized addition is rounded once and multiplication
by $2^{40}$ is exact. Its addition error is less than
$\varepsilon_{\mathrm{in}}/8$. Consequently
\[
 \varphi\le\frac M{2^{40}}
 \le\varphi+2\varepsilon_{\mathrm{ar}}+2\cdot2^{-40}.
\]
Since $v\le0$, the exact downward quotient used for
$Mv/2^{40}$ certifies a lower bound on the exact multiplier times
the required interpolant. This proves the rounding correctness of
both the leaf continuation and the subtree continuation bound used
for pruning.
